\chapter{2024年上海卷（秋）}

\section{填空题}

\begin{problem}
设全集 \(U=\{1,2,3,4,5\}\)，集合 \(A=\{2,4\}\)，则 \(\bar{A}=\)\fillinblank{}.

\end{problem}

\begin{answer}
\(\{1,3,5\}\)
\end{answer}

\begin{solution}
由题设，\(A\) 在全集 \(U\) 中的补集为
\(\bar{A}=\{1,3,5\}\).
故答案为 \(\{1,3,5\}\).
\end{solution}

\begin{problem}
已知 \(f(x)=\sqrt{x}\)（\(x>0\)），\(f(x)=1\)（\(x\le 0\)），则 \(f(3)=\)\fillinblank{}.

\end{problem}

\begin{answer}
\(\sqrt{3}\)
\end{answer}

\begin{solution}
因为 \(3>0\)，所以
\(f(3)=\sqrt{3}\).
故答案为 \(\sqrt{3}\).
\end{solution}

\begin{problem}
已知 \(x\in \R\)，则不等式 \(x^2-2x-3<0\) 的解集为\fillinblank{}.

\end{problem}

\begin{answer}
\(\{x\mid -1<x<3\}\)
\end{answer}

\begin{solution}
方程
\(x^2-2x-3=0\)
的两根为 \(-1\) 与 3，因此不等式 \(x^2-2x-3<0\) 的解集为
\(\{x\mid -1<x<3\}\).
\end{solution}

\begin{problem}
已知 \(f(x)=x^3+a\)，\(x\in \R\)，且 \(f(x)\) 是奇函数，则 \(a=\)\fillinblank{}.

\end{problem}

\begin{answer}
0
\end{answer}

\begin{solution}
因为 \(f(x)\) 为奇函数，所以
\(f(-x)+f(x)=0\).
即
\((-x)^3+a+x^3+a=0\)，
解得 \(a=0\).
\end{solution}

\begin{problem}
已知 \(k\in \R\)，\(\bs{a}=(2,5)\)，\(\bs{b}=(6,k)\)，且 \(\bs{a}//\bs{b}\)，则 \(k=\)\fillinblank{}.

\end{problem}

\begin{answer}
15
\end{answer}

\begin{solution}
因为 \(\bs{a}//\bs{b}\)，所以对应坐标成比例，即
\(2k=5\times 6\).
解得
\(k=15\).
\end{solution}

\begin{problem}
在 \((x+1)^n\) 的二项展开式中，若各项系数和为 32，则 \(x^2\) 项的系数为\fillinblank{}.

\end{problem}

\begin{answer}
10
\end{answer}

\begin{solution}
令 \(x=1\)，则
\((1+1)^n=32\)，
即
\(2^n=32\)，
解得 \(n=5\).

\((x+1)^5\) 的展开式通项为
\(T_{r+1}=\mathrm{C}_5^r x^{5-r}\).
令 \(5-r=2\)，得 \(r=3\)，故 \(x^2\) 项的系数为
\(\mathrm{C}_5^3=10\).
\end{solution}

\begin{problem}
已知抛物线 \(y^2=4x\) 上有一点 \(P\) 到准线的距离为 9，那么点 \(P\) 到 \(x\) 轴的距离为\fillinblank{}.

\end{problem}

\begin{answer}
\(4\sqrt{2}\)
\end{answer}

\begin{solution}
由抛物线 \(y^2=4x\) 的定义，准线方程为 \(x=-1\). 设点 \(P(x_0,y_0)\)，则
\(x_0+1=9\)，
解得 \(x_0=8\).

代入抛物线方程得
\(y_0^2=4x_0=32\)，
所以
\(y_0=\pm 4\sqrt{2}\).
因此点 \(P\) 到 \(x\) 轴的距离为 \(4\sqrt{2}\).
\end{solution}

\begin{problem}
某校举办科学竞技比赛，有 \(A,B,C\) 3 种题库，\(A\) 题库有 5000 道题，\(B\) 题库有 4000 道题，\(C\) 题库有 3000 道题. 小申已完成所有题，他 \(A\) 题库的正确率是 0.92，\(B\) 题库的正确率是 0.86，\(C\) 题库的正确率是 0.72. 现他从所有的题中随机选一题，正确率是\fillinblank{}.

\end{problem}

\begin{answer}
0.85
\end{answer}

\begin{solution}
根据题意，三种题库数量比为
\(5000:4000:3000=5:4:3\)，
故所占比重分别为
\(\frac{5}{12}\)，\(\frac{4}{12}\)，\(\frac{3}{12}\).
所以随机选一题答对的概率为
\(\frac{5}{12}\times 0.92+\frac{4}{12}\times 0.86+\frac{3}{12}\times 0.72=0.85\).
\end{solution}

\begin{problem}
已知虚数 \(z\)，其实部为 1，且 \(z+\frac{2}{z}=m\)（\(m\in \R\)），则实数 \(m=\)\fillinblank{}.

\end{problem}

\begin{answer}
2
\end{answer}

\begin{solution}
设
\(z=1+b\i\)，\(b\in \R,\ b\ne 0\).
则
\(z+\frac{2}{z}=1+b\i+\frac{2}{1+b\i} =\left( \frac{b^2+3}{1+b^2} \right)+\left( \frac{b^3-b}{1+b^2} \right)\i\).
因为 \(m\in \R\)，所以虚部为 0，即
\(\frac{b^3-b}{1+b^2}=0\).
解得 \(b=\pm 1\). 于是
\(m=\frac{b^2+3}{1+b^2}=2\).
\end{solution}

\begin{problem}
设集合 \(A\) 中的元素皆为无重复数字的三位正整数，且元素中任意两者之积皆为偶数，求集合中元素个数的最大值\fillinblank{}.

\end{problem}

\begin{answer}
329
\end{answer}

\begin{solution}
根据题意，集合中至多只有一个奇数，其余均是偶数.

首先讨论三位数中的偶数.

\circled{1} 当个位为 0 时，则百位和十位在剩余的 9 个数字中选择两个进行排列，则这样的偶数有
\(P_9^2=72\)
个；

\circled{2} 当个位不为 0 时，则个位有 \( \mathrm{C}_4^1 \) 个数字可选，百位有 \( \mathrm{C}_8^1 \) 个数字可选，十位有 \( \mathrm{C}_8^1 \) 个数字可选，

由分步乘法这样的偶数共有
\(\mathrm{C}_4^1\mathrm{C}_8^1\mathrm{C}_8^1=256\)
个.

最后再加上单独的奇数，所以集合中元素个数的最大值为
\(72+256+1=329\).
\end{solution}

\begin{problem}
已知点 \(B\) 在点 \(C\) 正北方向，点 \(D\) 在点 \(C\) 的正东方向，\(BC=CD\)，存在点 \(A\) 满足 \(\angle BAC=16.5^\circ\)，\(\angle DAC=37^\circ\)，则 \(\angle BCA=\)\fillinblank{}（精确到 0.1 度）.

\bitmapfigure[max width=0.42\linewidth,max totalheight=4.8cm]{img/2024/shanghai/q11_fig1.png}

\end{problem}

\begin{answer}
\(7.8^\circ\)
\end{answer}

\begin{solution}
设
\(\angle BCA=\theta\)，
则
\(\angle ACD=90^\circ-\theta\).

在 \(\triangle DCA\) 中，由正弦定理得
\(\frac{CA}{CD}=\frac{\sin 37^\circ}{\sin(90^\circ-\theta+37^\circ)}\).

在 \(\triangle BCA\) 中，由正弦定理得
\(\frac{CA}{CB}=\frac{\sin 16.5^\circ}{\sin(\theta+16.5^\circ)}\).

因为 \(CD=CB\)，所以
\(\frac{\sin 37^\circ}{\sin(90^\circ-\theta+37^\circ)} =\frac{\sin 16.5^\circ}{\sin(\theta+16.5^\circ)}\).
利用计算器解得
\(\theta\approx 7.8^\circ\).
故答案为 \(7.8^\circ\).
\end{solution}

\begin{problem}
无穷等比数列 \(\{a_n\}\) 满足首项 \(a_1>0\)，\(q>1\)，记 \(I_n=\{x-y\mid x,y\in [a_1,a_2]\cup [a_n,a_{n+1}]\}\)，若对任意正整数 \(n\) 集合 \(I_n\) 是闭区间，则 \(q\) 的取值范围是\fillinblank{}.

\end{problem}

\begin{answer}
\(q\ge 2\)
\end{answer}

\begin{solution}
由题设，
\(a_n=a_1q^{n-1}\).
因为 \(a_1>0\)，\(q>1\)，所以 \(a_{n+1}>a_n\).

当 \(n\ge 2\) 时，不妨设 \(x\ge y\). 则
\(x-y\in [0,a_2-a_1]\cup [a_n-a_2,a_{n+1}-a_1]\cup [0,a_{n+1}-a_n]\).
要使 \(I_n\) 为闭区间，就必须有
\(a_{n+1}-a_n\ge a_n-a_2\).
即
\(a_1q^{n-1}(q-2)+a_2\ge 0\).
等价于
\(q-2\ge -\frac{1}{q^{n-2}}\).
由于 \(q>1\)，当 \(n\to +\infty\) 时，\(-\frac{1}{q^{n-2}}\to 0\)，故必有
\(q-2\ge 0\).
因此
\(q\ge 2\).
\end{solution}
\section{单选题}

\begin{problem}
已知气候温度和海水表层温度相关，且相关系数为正数，对此描述正确的是（\quad）

\choices
    {气候温度高，海水表层温度就高}
    {气候温度高，海水表层温度就低}
    {随着气候温度由低到高，海水表层温度呈上升趋势}
    {随着气候温度由低到高，海水表层温度呈下降趋势}

\end{problem}

\begin{answer}
C
\end{answer}

\begin{solution}
A、B. 当气候温度高时，海水表层温度变高变低不确定，故 A、B 错误.

C、D. 因为相关系数为正，故随着气候温度由低到高时，海水表层温度呈上升趋势，故 C 正确，D 错误.

故选 C.
\end{solution}

\begin{problem}
下列函数 \(f(x)\) 的最小正周期是 \(2\pi\) 的是（\quad）

\choices
    {\(\sin x+\cos x\)}
    {\(\sin x\cos x\)}
    {\(\sin^2 x+\cos^2 x\)}
    {\(\sin^2 x-\cos^2 x\)}

\end{problem}

\begin{answer}
A
\end{answer}

\begin{solution}
因为
\(\sin x+\cos x=\sqrt{2}\sin\left( x+\frac{\pi}{4} \right)\)，
其最小正周期为 \(2\pi\).

又
\(\sin x\cos x=\frac{1}{2}\sin 2x\)，
最小正周期为 \(\pi\)；
\(\sin^2 x+\cos^2 x=1\)
为常值函数；
\(\sin^2 x-\cos^2 x=-\cos 2x\)
最小正周期为 \(\pi\). 故选 A.
\end{solution}

\begin{problem}
定义一个集合 \(\Omega\)，集合中的元素是空间内的点集，任取 \(P_1,P_2,P_3\in \Omega\)，存在不全为 0 的实数 \(\lambda_1,\lambda_2,\lambda_3\)，使得 \(\lambda_1\overrightarrow{OP_1}+\lambda_2\overrightarrow{OP_2}+\lambda_3\overrightarrow{OP_3}=\bs{0}\). 已知 \((1,0,0)\in \Omega\)，则 \((0,0,1)\notin \Omega\) 的充分条件是（\quad）

\choices
    {\((0,0,0)\in \Omega\)}
    {\((-1,0,0)\in \Omega\)}
    {\((0,1,0)\in \Omega\)}
    {\((0,0,-1)\in \Omega\)}

\end{problem}

\begin{answer}
C
\end{answer}

\begin{solution}
根据题意，这三个向量 \(\overrightarrow{OP_1}\)，\(\overrightarrow{OP_2}\)，\(\overrightarrow{OP_3}\) 共面，即这三个向量不能构成空间的一个基底.

A. 在空间直角坐标系中，\((0,0,0)\)、\((1,0,0)\)、\((0,0,1)\) 三个向量共面，则当 \((-1,0,0)\)、\((1,0,0)\in \Omega\) 时无法推出 \((0,0,1)\notin \Omega\)，A 错误；

B. 在空间直角坐标系中，\((-1,0,0)\)、\((1,0,0)\)、\((0,0,1)\) 三个向量共面，则当 \((0,0,0)\)、\((1,0,0)\in \Omega\) 时无法推出 \((0,0,1)\notin \Omega\)，B 错误；

C. 在空间直角坐标系中，\((1,0,0)\)、\((0,0,1)\)、\((0,1,0)\) 三个向量不共面，可构成空间的一个基底，则由 \((1,0,0)\)、\((0,1,0)\in \Omega\) 能推出 \((0,0,1)\notin \Omega\)，C 正确；

D. 在空间直角坐标系中，\((1,0,0)\)、\((0,0,1)\)、\((0,0,-1)\) 三个向量共面，则当 \((0,0,-1)\)、\((1,0,0)\in \Omega\) 时无法推出 \((0,0,1)\notin \Omega\)，D 错误.

故选 C.
\end{solution}

\begin{problem}
已知函数 \(f(x)\) 的定义域为 \(\R\)，定义集合 \(M=\{x_0\mid x_0\in \R,\ \forall x\in(-\infty,x_0),\ f(x)<f(x_0)\}\). 在使得 \(M=[-1,1]\) 的所有 \(f(x)\) 中，下列成立的是（\quad）

\choices
    {存在 \(f(x)\) 是偶函数}
    {存在 \(f(x)\) 在 \(x=2\) 处取最大值}
    {存在 \(f(x)\) 是严格增函数}
    {存在 \(f(x)\) 在 \(x=-1\) 处取到极小值}

\end{problem}

\begin{answer}
B
\end{answer}

\begin{solution}
A：若 \(f(x)\) 是偶函数，取 \(x_0=1\in[-1,1]\)，则对任意 \(x<1\) 有 \(f(x)<f(1)\)，但 \(f(-1)=f(1)\)，矛盾.

B：可构造
\[
f(x)=
\begin{cases}
-2, & x<-1,\\
x, & -1\le x\le 1,\\
1, & x>1,
\end{cases}
\]
此时 \(M=[-1,1]\)，且 \(f(2)=1\) 为最大值.

C：若 \(f(x)\) 严格递增，则对任意 \(x_0\in \R\) 均满足定义，因此 \(M=\R\)，矛盾.

D：若 \(f(x)\) 在 \(x=-1\) 处取极小值，则在 \(-1\) 左邻域内存在 \(x\) 使 \(f(x)>f(-1)\)，与 \(M\) 的定义矛盾.

故选 B.
\end{solution}
\section{解答题}

\begin{problem}
如图为正四棱锥 \(P-ABCD\)，\(O\) 为底面 \(ABCD\) 的中心.

\begin{enumerate}
\item 若 \(AP=5\)，\(AD=3\sqrt{2}\)，求 \(\triangle POA\) 绕 \(PO\) 旋转一周形成的几何体的体积；
\item 若 \(AP=AD\)，\(E\) 为 \(PB\) 的中点，求直线 \(BD\) 与平面 \(AEC\) 所成角的大小.
\end{enumerate}

\bitmapfigure[max width=0.55\linewidth,max totalheight=4.8cm]{img/2024/shanghai/q17_fig1.png}

\end{problem}

\begin{solution}
（1）因为正四棱锥满足 \(PO\perp \text{平面 }ABCD\)，而 \(AO\subset \text{平面 }ABCD\)，故
\(PO\perp AO\).
又底面 \(ABCD\) 是正方形，且 \(AD=3\sqrt{2}\)，所以
\(AO=3\).
由勾股定理，
\(PO=\sqrt{PA^2-AO^2}=\sqrt{25-9}=4\).
\(\triangle POA\) 绕 \(PO\) 旋转一周形成的是圆锥，其高为 4，底面半径为 3，因此体积为
\(V=\frac{1}{3}\pi\times 3^2\times 4=12\pi\).

（2）连接 \(EA,EO,EC\). 根据题意结合正四棱锥的性质，每个侧面都是等边三角形.

\bitmapfigure[max width=0.55\linewidth,max totalheight=4.8cm]{img/2024/shanghai/q17_solution_fig1.png}

由 \(E\) 是 \(PB\) 的中点，得
\(AE\perp PB\)，\(CE\perp PB\).
又 \(AE\cap CE=E\)，且 \(AE,CE\subset \text{平面 }ACE\)，故
\(PB\perp \text{平面 }ACE\).
于是 \(BE\perp \text{平面 }ACE\). 又
\(BD\cap \text{平面 }ACE=O\)，
所以直线 \(BD\) 与平面 \(AEC\) 所成角为 \(\angle BOE\).

不妨设
\(AP=AD=6\)，
则
\(BO=3\sqrt{2}\)，\(BE=3\).
因此
\(\sin \angle BOE=\frac{BE}{BO}=\frac{3}{3\sqrt{2}}=\frac{\sqrt{2}}{2}\).
又线面角属于 \(\left[ 0,\frac{\pi}{2} \right]\)，故
\(\angle BOE=\frac{\pi}{4}\).
\end{solution}

\begin{problem}
若 \(f(x)=\log_a x\)（\(a>0,\ a\ne 1\)）.

\begin{enumerate}
\item \(y=f(x)\) 过 \((4,2)\)，求 \(f(2x-2)<f(x)\) 的解集；
\item 存在 \(x\) 使得 \(f(x+1)\)、\(f(ax)\)、\(f(x+2)\) 成等差数列，求 \(a\) 的取值范围.
\end{enumerate}

\end{problem}

\begin{solution}
（1）因为 \(y=f(x)\) 过 \((4,2)\)，所以
\(\log_a 4=2\)，
即
\(a^2=4\).
又 \(a>0,\ a\ne 1\)，故 \(a=2\).

于是 \(f(x)=\log_2 x\) 在 \((0,+\infty)\) 上单调递增，所以
\(f(2x-2)<f(x) \iff 0<2x-2<x\).
解得
\(1<x<2\).

（2）由 \(f(x+1)\)、\(f(ax)\)、\(f(x+2)\) 成等差数列，得
\(2f(ax)=f(x+1)+f(x+2)\).
即
\(2\log_a(ax)=\log_a(x+1)+\log_a(x+2)\).
因此
\(a^2x^2=(x+1)(x+2)\).
又 \(x>0\)，所以
\[
\begin{aligned}
a^2
&=\frac{x^2+3x+2}{x^2} \\
&=1+\frac{3}{x}+\frac{2}{x^2} \\
&=2\left( \frac{1}{x}+\frac{3}{4} \right)^2-\frac{1}{8}
\end{aligned}
\]
令
\(t=\frac{1}{x}\in(0,+\infty)\)，
则
\(a^2=2\left( t+\frac{3}{4} \right)^2-\frac{1}{8}\).
其在 \((0,+\infty)\) 上的值域为 \((1,+\infty)\)，故
\(a^2>1\).
由于 \(a>0\)，得
\(a>1\).
\end{solution}

\begin{problem}
为了解某地初中学生体育锻炼时长与学业成绩的关系，从该地区 29000 名学生中抽取 580 人，得到日均体育锻炼时长与学业成绩的数据如下表所示：

\begin{center}
\begin{tabular}{|c|c|c|c|c|c|}
\hline
时间范围/学业成绩 & \([0,0.5)\) & \([0.5,1)\) & \([1,1.5)\) & \([1.5,2)\) & \([2,2.5)\) \\
\hline
优秀 & 5 & 44 & 42 & 3 & 1 \\
\hline
不优秀 & 134 & 147 & 137 & 40 & 27 \\
\hline
\end{tabular}
\end{center}

\begin{enumerate}
\item 该地区 29000 名学生中体育锻炼时长不少于 1 小时人数约为多少？
\item 估计该地区初中学生日均体育锻炼的时长（精确到 0.1）；
\item 是否有 95\% 的把握认为学业成绩优秀与日均体育锻炼时长不小于 1 小时且小于 2 小时有关？
（附：\(\chi^2=\frac{n(ad-bc)^2}{(a+b)(c+d)(a+c)(b+d)}\)，其中 \(n=a+b+c+d\)，\(P(\chi^2\ge 3.841)\approx 0.05\).）
\end{enumerate}

\end{problem}

\begin{solution}
（1）锻炼时长不少于 1 小时的人数占比为
\(\frac{179+43+28}{580}=\frac{25}{58}\).
故估计该地区 29000 名学生中体育锻炼时长不少于 1 小时的人数约为
\(29000\times \frac{25}{58}=12500\).

（2）估计日均体育锻炼时长约为
\[
\frac{1}{580}\left[ \frac{0.5}{2}\times 139 +\frac{0.5+1}{2}\times 191 +\frac{1+1.5}{2}\times 179 +\frac{1.5+2}{2}\times 43 +\frac{2+2.5}{2}\times 28 \right]\approx 0.9
\]
所以估计该地区初中学生日均体育锻炼的时长为 0.9 小时.

（3）由题列联表如下：

\begin{center}
\begin{tabular}{|c|c|c|c|}
\hline
 & \([1,2)\) & 其他 & 合计 \\
\hline
优秀 & 45 & 50 & 95 \\
\hline
不优秀 & 177 & 308 & 485 \\
\hline
合计 & 222 & 358 & 580 \\
\hline
\end{tabular}
\end{center}

提出零假设 \(H_0\)：该地区成绩优秀与日均锻炼时长不少于 1 小时但少于 2 小时无关，其中 \(\alpha=0.05\).

\[
\chi^2=\frac{580\times (45\times 308-177\times 50)^2}{95\times 485\times 222\times 358}\approx 3.976>3.841
\]
故零假设不成立，即有 95\% 的把握认为学业成绩优秀与日均体育锻炼时长不小于 1 小时且小于 2 小时有关.
\end{solution}

\begin{problem}
已知双曲线 \(\Gamma:x^2-\frac{y^2}{b^2}=1\)（\(b>0\)），左右顶点分别为 \(A_1,A_2\)，过点 \(M(-2,0)\) 的直线 \(l\) 交双曲线 \(\Gamma\) 于 \(P,Q\) 两点.

\begin{enumerate}
\item 若离心率 \(e=2\) 时，求 \(b\) 的值；
\item 若 \(b=\frac{2\sqrt{6}}{3}\)，\(\triangle MA_2P\) 为等腰三角形时，且点 \(P\) 在第一象限，求点 \(P\) 的坐标；
\item 连接 \(OQ\) 并延长，交双曲线 \(\Gamma\) 于点 \(R\)，若 \(\overrightarrow{A_1R}\cdot\overrightarrow{A_2P}=1\)，求 \(b\) 的取值范围.
\end{enumerate}

\end{problem}

\begin{solution}
（1）对于双曲线
\(x^2-\frac{y^2}{b^2}=1\)，
有 \(a=1\). 因为 \(e=\frac{c}{a}=2\)，所以 \(c=2\). 又 \(c^2=a^2+b^2\)，故
\[
b=\sqrt{c^2-a^2}=\sqrt{4-1}=\sqrt{3}
\]

（2）当 \(b=\frac{2\sqrt{6}}{3}\) 时，双曲线方程为
\(x^2-\frac{3y^2}{8}=1\).
其中 \(M(-2,0)\)，\(A_2(1,0)\).

若以 \(MA_2\) 为底，则点 \(P\) 落在直线 \(x=-\frac{1}{2}\) 上，与点 \(P\) 在第一象限矛盾，舍去.

若以 \(A_2P\) 为底，则 \(MP=MA_2=3\). 联立
\[
\begin{cases}
x^2-\frac{3y^2}{8}=1,\\
(x+2)^2+y^2=9,
\end{cases}
\]
解得
\(\left( -\frac{23}{11},\frac{8\sqrt{17}}{11} \right)\)、\(\left(-\frac{23}{11},-\frac{8\sqrt{17}}{11}\right)\)、\((1,0)\).
这些点都不在第一象限，故舍去.

若以 \(MP\) 为底，则 \(A_2P=MA_2=3\). 设 \(P(x_0,y_0)\)，其中 \(x_0>0\)，\(y_0>0\)，则
\[
\begin{cases}
(x_0-1)^2+y_0^2=9,\\
x_0^2-\frac{3y_0^2}{8}=1.
\end{cases}
\]
解得
\(x_0=2\)，\(y_0=2\sqrt{2}\).
故
\(P(2,2\sqrt{2})\).

（3）根据题意，\(A_1(-1,0)\)，\(A_2(1,0)\). 当直线 \(l\) 的斜率为 0 时，不满足 \(\overrightarrow{A_1R}\cdot\overrightarrow{A_2P}=1\)，故设
\(l:x=my-2\).
设
\(P(x_1,y_1)\)，\(Q(x_2,y_2)\).
由双曲线的对称性知
\(R(-x_2,-y_2)\).

联立
\[
\begin{cases}
x=my-2,\\
x^2-\frac{y^2}{b^2}=1
\end{cases}
\]
可得关于 \(y\) 的方程
\((b^2m^2-1)y^2-4b^2my+3b^2=0\).
因此
\(y_1+y_2=\frac{4b^2m}{b^2m^2-1}\)，\(y_1y_2=\frac{3b^2}{b^2m^2-1}\).

\bitmapfigure[max width=0.55\linewidth,max totalheight=4.8cm]{img/2024/shanghai/q20_solution_fig1.png}

又
\(\overrightarrow{A_1R}=(-x_2+1,-y_2)\)，\(\overrightarrow{A_2P}=(x_1-1,y_1)\)，
由 \(\overrightarrow{A_1R}\cdot\overrightarrow{A_2P}=1\)，并结合 \(x_1=my_1-2\)，\(x_2=my_2-2\)，得
\[
\begin{aligned}
&(-x_2+1)(x_1-1)-y_1y_2=1 \\
\iff{}& -(my_2-3)(my_1-3)-y_1y_2=1 \\
\iff{}& y_1y_2(m^2+1)-3m(y_1+y_2)+10=0
\end{aligned}
\]
将 \(y_1+y_2=\frac{4b^2m}{b^2m^2-1}\)，\(y_1y_2=\frac{3b^2}{b^2m^2-1}\) 代入上式，得
\[
\begin{aligned}
&(m^2+1)\cdot \frac{3b^2}{b^2m^2-1}-3m\cdot \frac{4b^2m}{b^2m^2-1}+10=0 \\
\iff{}& b^2m^2+3b^2-10=0
\end{aligned}
\]
所以
\(m^2=\frac{10}{b^2}-3\).
由 \(m^2\ge 0\) 得
\(0<b^2\le \frac{10}{3}\).
又由 \(b^2m^2-1\ne 0\) 得
\(10-3b^2\ne 1\)，
即
\(b^2\ne 3\).
故
\(b^2\in \left( 0,3 \right)\cup \left( 3,\frac{10}{3} \right]\).
因为 \(b>0\)，所以
\(b\in (0,\sqrt{3})\cup \left( \sqrt{3},\frac{\sqrt{30}}{3} \right]\).
\end{solution}

\begin{problem}
对于一个函数 \(f(x)\) 和一个点 \(M(a,b)\)，令 \(s(x)=(x-a)^2+\left( f(x)-b \right)^2\). 若 \(P(x_0,f(x_0))\) 是 \(s(x)\) 取到最小值的点，则称 \(P\) 是 \(M\) 在 \(f(x)\) 上的“最近点”.

\begin{enumerate}
\item 对于 \(f(x)=\frac{1}{x}\)（\(x>0\)），求证：对于点 \(M(0,0)\)，存在点 \(P\)，使得点 \(P\) 是 \(M\) 在 \(f(x)\) 上的“最近点”；
\item 对于 \(f(x)=\e^x\)，\(M(1,0)\)，请判断是否存在一个点 \(P\)，它是 \(M\) 在 \(f(x)\) 上的“最近点”，且直线 \(MP\) 与 \(y=f(x)\) 在点 \(P\) 处的切线垂直；
\item 已知 \(y=f(x)\) 在定义域 \(\R\) 上存在导函数 \(f'(x)\)，且函数 \(g(x)\) 在定义域 \(\R\) 上恒正，设点 \(M_1(t-1,f(t)-g(t))\)，\(M_2(t+1,f(t)+g(t))\). 若对任意的 \(t\in \R\)，存在点 \(P\) 同时是 \(M_1,M_2\) 在 \(f(x)\) 上的“最近点”，试判断 \(f(x)\) 的单调性.
\end{enumerate}

\end{problem}

\begin{solution}
（1）当 \(M(0,0)\) 时，
\[
s(x)=x^2+\left( \frac{1}{x} \right)^2=x^2+\frac{1}{x^2}\ge 2\sqrt{x^2\cdot \frac{1}{x^2}}=2
\]
当且仅当
\(x^2=\frac{1}{x^2}\)，
即 \(x=1\) 时取等号. 因此存在点
\(P(1,1)\)，
使它是 \(M(0,0)\) 在 \(f(x)\) 上的“最近点”.

（2）由题设，
\(s(x)=(x-1)^2+\left( \e^x-0 \right)^2=(x-1)^2+\e^{2x}\).
则
\(s'(x)=2(x-1)+2\e^{2x}\).
因为 \(2(x-1)\) 与 \(2\e^{2x}\) 都是严格增函数，所以 \(s'(x)\) 也是严格增函数. 又
\(s'(0)=0\)，
因此 \(x=0\) 为 \(s(x)\) 的最小值点，此时
\(P(0,1)\).

又
\(f'(x)=\e^x\)，
故 \(f(x)\) 在点 \(P\) 处的切线斜率为
\(f'(0)=1\).
而
\(k_{MP}=\frac{0-1}{1-0}=-1\)，
所以
\(k_{MP}\cdot 1=-1\).
因此直线 \(MP\) 与切线垂直. 故这样的点存在，且
\(P(0,1)\).

（3）设
\(s_1(x)=(x-t+1)^2+\left( f(x)-f(t)+g(t) \right)^2\)，
\(s_2(x)=(x-t-1)^2+\left( f(x)-f(t)-g(t) \right)^2\).
若点 \(P(x_0,f(x_0))\) 同时是 \(M_1,M_2\) 的最近点，则 \(x_0\) 同时是 \(s_1,s_2\) 的最小值点，所以
\(s_1'(x_0)=0\)，\(s_2'(x_0)=0\).
即
\[
\begin{cases}
2(x_0-t+1)+2f'(x_0)\left( f(x_0)-f(t)+g(t) \right)=0,\\
2(x_0-t-1)+2f'(x_0)\left( f(x_0)-f(t)-g(t) \right)=0.
\end{cases}
\]
两式相减得
\(1+f'(x_0)g(t)=0\)，
故
\(f'(x_0)=-\frac{1}{g(t)}<0\).

再由“最近点”定义，有
\(s_1(x_0)\le s_1(t)\)，\(s_2(x_0)\le s_2(t)\).
两式相加可得
\((x_0-t)^2+\left( f(x_0)-f(t) \right)^2\le 0\).
因此
\(x_0=t\)，\(f(x_0)=f(t)\).
于是
\(f'(t)=f'(x_0)=-\frac{1}{g(t)}<0\).
由于 \(t\in \R\) 任意，故
\(f'(t)<0\)，\(\text{对任意 }t\in \R\text{ 都成立}\).
所以 \(f(x)\) 在 \(\R\) 上严格单调递减.
\end{solution}
